% TOP_PROBLEM: {"id": 30000395, "title": "Hoop Conjecture for Black-Hole Formation", "category_id": 16, "category": {"id": 16, "name": "physics", "display_name": "Mathematical Physics", "description": "Problems at the intersection of mathematics and physics.", "slug": "physics", "order_index": 16, "created_at": "2026-07-31T15:26:25.671Z"}, "status": "open", "statusQualification": "The primary report explicitly says that mass and circumference are not clearly defined in this formulation. Event and apparent horizons are different notions. No arbitrary definition of a hoop, quasi-local mass, or horizon is imposed here, and the time-symmetric results do not resolve every interpretation."}
% FIRST_POSTED: 2026-10-11
% ENTRY_KIND: statement_only
% SUBMISSION: {"kind": "statement_only", "category": "top_problems", "problem_id": "30000395", "proposal_type": "established", "source": "Mathematical Aspects of General Relativity (2006), Oberwolfach Reports, DOI 10.4171/owr/2006/02, pp. 73--118; official report source and mathematical context. https://doi.org/10.4171/owr/2006/02; Official EMS Press PDF of the same Oberwolfach report; PDF pages 16--17, printed pages 88--89. https://ems.press/content/serial-article-files/46032", "status_evidence": "Saved source status is open; retained unchanged. The official historical OWR source was inspected on 2026-10-11. The saved literature assessment is dated 2026-08-21 and is not a fresh comprehensive status verification. ", "status_checked_on": "2026-10-11", "duplicate_check": "Checked canonical identity and normalized formulations against the complete saved UnsolvedMath registry; No exact normalized duplicate found. Mathematical equivalence and current literature were not exhaustively searched; no new ID allocated.", "attribution": "Kip Thorne (1972), as discussed by Yuguang Shi with Luen-Fai Tam. Statement preparation and source audit assisted by OpenAI GPT-6 (Codex); no human mathematical authorship or independent proof verification claimed."}
% SOURCE_STATEMENT: "Do black holes with horizons form if and only if a mass $M$ is compacted into a region whose circumference in every direction is at most $4\\pi M$?"
% Preparation and final self-audit: CONTRIBUTING.md prompts C and D; source/audit artifacts retained in the local physics_statements_2026_10_11 batch.
% Imported from: physics_statements.tex; UnsolvedMath ID 30000395
% !TeX program = lualatex
\documentclass[11pt]{article}
\usepackage[margin=25mm]{geometry}
\usepackage{fontspec}
\setmainfont{Latin Modern Roman}
\usepackage{amsmath,amssymb,mathtools}
\usepackage{unicode-math}
\setmathfont{Latin Modern Math}
\usepackage{xurl,hyperref}
\hypersetup{colorlinks=true,urlcolor=blue}
\setcounter{secnumdepth}{0}
\setlength{\parindent}{0pt}
\setlength{\parskip}{5pt}
\setlength{\emergencystretch}{3em}
\newcommand{\sref}[2]{\hyperlink{source:#1:#2}{[S#2]}}
\newcommand{\cataloguescope}[1]{\paragraph{Recorded target.}#1}
\newcommand{\R}{\mathbb R}
\newcommand{\C}{\mathbb C}
\newcommand{\N}{\mathbb N}
\newcommand{\Z}{\mathbb Z}
\newcommand{\Q}{\mathbb Q}
\newcommand{\F}{\mathbb F}
\newcommand{\eps}{\varepsilon}
\begin{document}
% UnsolvedMath ID 30000395; source catalogue code OWR-1184-001
\setcounter{section}{30000394}
\section{Hoop Conjecture for Black-Hole Formation}\label{30000395}
{\small UnsolvedMath ID 30000395 · Mathematical Physics}
\subsection{Definitions and mathematical statement}

\paragraph{Shared notation.}$\mathbb N=\{1,2,\ldots\}$; the standard number fields have their usual meanings. All additional conventions are stated locally.


\paragraph{Definitions and conventions.} Use geometrized units $G=c=1$, so a mass $M$ has units of length. An event horizon is the boundary separating points that can send causal signals to future null infinity from the black-hole region. The hoop conjecture proposes a relation between compactness, a notion of mass, and circumferences in all directions. In time-symmetric initial data an apparent horizon corresponds to an appropriate closed minimal surface, which is a more specific geometric setting discussed in the report.

\paragraph{Source scope and qualifications.} The primary report explicitly says that mass and circumference are not clearly defined in this formulation. Event and apparent horizons are different notions. No arbitrary definition of a hoop, quasi-local mass, or horizon is imposed here, and the time-symmetric results do not resolve every interpretation.

\paragraph{Catalogue mathematical target.}
\paragraph{Catalogue wording retained for traceability.}
Do black holes with horizons form if and only if a mass $M$ is compacted into a region whose circumference in every direction is at most $4\pi M$?

\paragraph{Saved literature context.} The supplied catalogue's assessment dated 2026-08-21 records: The extracted item is explicitly labelled a loosely formulated conjecture. Without the omitted mathematical formulation there is no verified theorem or counterexample to classify it as solved or disproved. The assessment is retained as dated provenance; it is not a new comprehensive literature or open-status certification.
\subsection{Short English statement}
Make precise and test the claim that horizons form exactly when mass is compressed inside sufficiently short hoops in every direction.
\subsection{Sources}
\begin{itemize}
\item[\hypertarget{source:30000395:1}{S1}] ULAM.AI and the UnsolvedMath Contributors, UnsolvedMath, record 30000395 (OWR-1184-001). Catalogue curation under CC BY 4.0; definitions and English exposition added locally.
\url{https://huggingface.co/datasets/ulamai/UnsolvedMath}.
\item[\hypertarget{source:30000395:2}{S2}] Mathematical Aspects of General Relativity (2006), Oberwolfach Reports, DOI 10.4171/owr/2006/02, pp. 73--118; official report source and mathematical context.
\url{https://doi.org/10.4171/owr/2006/02}.
\item[\hypertarget{source:30000395:3}{S3}] Official EMS Press PDF of the same Oberwolfach report; PDF pages 16--17, printed pages 88--89.
\url{https://ems.press/content/serial-article-files/46032}.
\end{itemize}
\label{end:30000395}
\end{document}
